% !TeX program = pdflatex
% !TeX encoding = UTF-8
\documentclass[11pt,aspectratio=169,t]{beamer}
% Все презентации курса компилируются только pdflatex.
\usepackage{iftex}
\RequirePDFTeX
\usepackage{cmap}
\usepackage[T2A]{fontenc}
\usepackage[utf8]{inputenc}
\usepackage[english,russian]{babel}
\usepackage{PTSerif}
\usepackage{amsmath,amsfonts,amssymb}
\usepackage{graphicx,microtype,anyfontsize,tikz}
\usepackage{ragged2e}
\usetikzlibrary{arrows.meta,calc}

\usetheme{default}
\usefonttheme{serif}
\DeclareMathOperator{\sen}{sen}
\bibliographystyle{apalike}

% Общие сведения; номер лекции, её название и дата задаются в N.tex.
\newcommand{\coursename}{Введение в маневрирование КА}
\author{Чиняев Владимир Николаевич}
\institute{Московский физико-технический институт}

% Книжное оформление с голубым акцентом.
\definecolor{coursepaper}{HTML}{FBF9F3}
\definecolor{courseink}{HTML}{302B28}
\definecolor{courseaccent}{HTML}{187AA5}
\definecolor{coursemuted}{HTML}{6D6560}
\definecolor{courserule}{HTML}{A9CADA}
\definecolor{coursetint}{HTML}{EDF4F6}
\tikzset{
  course diagram/.style={font=\fontsize{9.5}{11.5}\selectfont,
    line cap=round,line join=round,>=Latex},
  orbit line/.style={draw=courseaccent,line width=1pt},
  aux line/.style={draw=courserule,line width=.5pt},
  vector line/.style={draw=courseink,line width=.65pt,->},
  angle line/.style={draw=courseaccent,line width=.8pt,->},
  orbit point/.style={circle,fill=courseaccent,inner sep=1.7pt}
}
\setbeamercolor{normal text}{fg=courseink,bg=coursepaper}
\setbeamercolor{background canvas}{bg=coursepaper}
\setbeamercolor{structure}{fg=courseaccent}
\setbeamercolor{alerted text}{fg=courseaccent}
\setbeamercolor{frametitle}{fg=courseink,bg=coursepaper}
\setbeamercolor{footline}{fg=coursemuted,bg=coursepaper}
\setbeamercolor{caption name}{fg=courseaccent}

% Основной текст — 11/14 pt; подписи к рисункам — 9.5/11.5 pt.
\let\coursenormalsize\normalsize
\renewcommand{\normalsize}{\coursenormalsize\fontsize{11}{14}\selectfont}
\setbeamerfont{normal text}{size*={11pt}{14pt}}
\setbeamerfont{frametitle}{size*={20pt}{23pt},series=\mdseries}
\setbeamerfont{title}{size*={27pt}{31pt},series=\mdseries}
\setbeamerfont{subtitle}{size*={12.8pt}{15pt},series=\mdseries}
\setbeamerfont{author}{size*={10.8pt}{13pt}}
\setbeamerfont{institute}{size*={9.2pt}{11pt}}
\setbeamerfont{date}{size*={9pt}{11pt}}
\setbeamerfont{footline}{size*={6.8pt}{8pt}}
\setbeamerfont{caption}{size*={9.5pt}{11.5pt}}
\setbeamerfont{itemize/enumerate body}{size*={11pt}{14pt}}
\setbeamerfont{itemize/enumerate subbody}{size*={11pt}{14pt}}
\setbeamerfont{itemize/enumerate subsubbody}{size*={11pt}{14pt}}
\AtBeginDocument{\normalsize}

\setbeamersize{text margin left=6mm,text margin right=6mm}
\setbeamertemplate{navigation symbols}{}
\setbeamertemplate{headline}{}
\setbeamertemplate{caption}[numbered]
\setbeamertemplate{itemize items}[circle]

% Стандартные блоки Beamer позволяют группировать текст и формулы.
\setbeamertemplate{blocks}[default]
\setbeamerfont{block title}{size*={11pt}{14pt},series=\bfseries}
\setbeamerfont{block body}{size*={11pt}{14pt}}
\setbeamercolor{block title}{fg=courseaccent,bg=coursepaper}
\setbeamercolor{block body}{fg=courseink,bg=coursetint}
\setbeamercolor{block title alerted}{parent=block title}
\setbeamercolor{block body alerted}{parent=block body}
\setbeamercolor{block title example}{parent=block title}
\setbeamercolor{block body example}{parent=block body}

% Компактный заголовок; его высота учитывается при размещении содержания.
\setbeamertemplate{frametitle}{%
  \vskip4.5mm%
  \begin{beamercolorbox}[wd=\textwidth,sep=0pt]{frametitle}
    \usebeamerfont{frametitle}\strut\insertframetitle\strut\par
    \vskip2mm%
    {\color{courseaccent}\hrule height .55pt}
  \end{beamercolorbox}%
}

% Нижний колонтитул: курс слева, лекция и номер слайда справа.
\setbeamertemplate{footline}{%
  \begin{beamercolorbox}[wd=\paperwidth,ht=5.5mm,dp=0pt]{footline}
    \begin{tikzpicture}[remember picture,overlay,x=1mm,y=1mm]
      \begin{scope}[shift={(current page.south west)}]
        \draw[courserule,line width=.3pt] (6,5.5)--(154,5.5);
        \node[text=coursemuted,anchor=west,inner sep=0pt,
          font=\usebeamerfont{footline}] at (6,2.8) {\coursename};
        \node[text=coursemuted,anchor=east,inner sep=0pt,
          font=\usebeamerfont{footline}] at (154,2.8)
          {Лекция \arabic{lecture}\quad\textbullet\quad\insertframenumber};
      \end{scope}
    \end{tikzpicture}%
  \end{beamercolorbox}%
}

% Титульник рассчитан на 16:9 (160 × 90 мм), без нижнего колонтитула.
\setbeamertemplate{title page}{%
  \begin{tikzpicture}[remember picture,overlay,x=1mm,y=1mm]
    \begin{scope}[shift={(current page.south west)}]
      \node[text=coursemuted,anchor=north,font=\usebeamerfont{institute}]
        at (80,82) {\insertinstitute};
      \draw[courseaccent,line width=.55pt] (71,71)--(89,71);
      \node[text=courseink,align=center,text width=125mm,
        font=\usebeamerfont{title}] at (80,53)
        {\hyphenpenalty=10000\coursename\par};
      \node[text=courseaccent,align=center,text width=145mm,font=\usebeamerfont{subtitle}]
        at (80,32) {Лекция \arabic{lecture}. \inserttitle};
      \node[text=courseink,font=\usebeamerfont{author}]
        at (80,19) {\insertauthor};
      \node[text=coursemuted,font=\usebeamerfont{date}]
        at (80,9) {\insertdate};
    \end{scope}
  \end{tikzpicture}%
}

\usepackage{booktabs,array}
\setlength{\abovedisplayskip}{5pt}
\setlength{\belowdisplayskip}{5pt}
\setlength{\abovedisplayshortskip}{3pt}
\setlength{\belowdisplayshortskip}{4pt}
\newcommand{\eqcompact}{\fontsize{10}{12}\selectfont}
\newcommand{\figcaption}[1]{{\usebeamerfont{caption}\color{coursemuted}#1\par}}
\newcommand{\dvr}{\Delta\widehat V_r}
\newcommand{\dvt}{\Delta\widehat V_t}
\newcommand{\drh}{\Delta\widehat r}
\newcommand{\vz}{\widehat V_z}
\newcommand{\zh}{\widehat z}
\newcommand{\imp}{\text{\normalfont имп}}
\setcounter{lecture}{4}
\title{Околокруговое приближение.\\Уравнения движения и импульсный переход}
\date{13 октября 2026 г.}

\begin{document}

% Только интервалы около формул; основной текст остаётся 11/14 pt.
\newcommand{\derivationmathspacing}{%
  \setlength{\abovedisplayskip}{2pt}%
  \setlength{\belowdisplayskip}{3pt}%
  \setlength{\abovedisplayshortskip}{2pt}%
  \setlength{\belowdisplayshortskip}{2pt}%
  \setlength{\jot}{2pt}}

\begin{frame}[plain]
\titlepage
\end{frame}

\begin{frame}{Модель движения}
\relax
\justifying
Пусть $\boldsymbol R$ --- радиус-вектор аппарата относительно центра
притяжения, $\boldsymbol V$ --- его скорость в инерциальной системе
координат, $\boldsymbol a_{\text{упр}}$ --- управляющее ускорение.
Исходные уравнения имеют вид
\par\medskip
\derivationmathspacing
\[
\boxed{\begin{aligned}
 \dot{\boldsymbol R}&=\boldsymbol V,\\[5pt]
 \dot{\boldsymbol V}&=-\frac{\mu}{\|\boldsymbol R\|^3}\boldsymbol R
                  +\boldsymbol a_{\text{упр}}.
\end{aligned}}
\]
\par\medskip
Точкой обозначается производная по физическому времени $t$.
\end{frame}

\begin{frame}{Координаты и локальный базис}
\relax
\begin{columns}[T,onlytextwidth]
\begin{column}{.61\textwidth}
\centering
\resizebox{\linewidth}{!}{% Рисунок дословно из manual/near-circular-transfer.tex.
\begingroup
\definecolor{figink}{HTML}{25323F}
\definecolor{figgray}{HTML}{687582}
\definecolor{figrad}{HTML}{176AA6}
\definecolor{figtan}{HTML}{087F78}
\definecolor{fignormal}{HTML}{B57714}
\definecolor{figplane}{HTML}{EEF4F8}
\definecolor{figorbit}{HTML}{79539C}
\begin{tikzpicture}[x=1cm,y=1cm,font=\small,
 line cap=round,line join=round,
 fixed/.style={draw=figgray,line width=.65pt,-{Stealth[length=2.1mm]}},
 local/.style={line width=1.35pt,-{Stealth[length=2.7mm,width=1.8mm]}},
 guide/.style={draw=figgray!60,line width=.6pt,dash pattern=on 2.4pt off 2.4pt},
 point/.style={circle,inner sep=2pt,fill=figink},
 lab/.style={fill=white,inner sep=1.5pt}]
\path[use as bounding box] (0,.1) rectangle (15.9,10.45);
\pgfmathsetmacro{\orba}{4.3}
\pgfmathsetmacro{\orbi}{30}
\pgfmathsetmacro{\orbv}{30}
\pgfmathsetmacro{\px}{\orba*cos(\orbv)}
\pgfmathsetmacro{\py}{\orba*sin(\orbv)*cos(\orbi)}
\pgfmathsetmacro{\pz}{\orba*sin(\orbv)*sin(\orbi)}
\pgfmathsetmacro{\phiu}{atan2(\py,\px)}
\begin{scope}[shift={(7.5,5.2)},
 x={(1cm,-.16cm)},y={(-.5cm,.58cm)},z={(0cm,1cm)}]
 \coordinate (O) at (0,0,0);
 \coordinate (H) at (\px,\py,0);
 \coordinate (P) at (\px,\py,\pz);
 \coordinate (Q) at ({\orba*cos(50)},{\orba*sin(50)},0);
 \filldraw[fill=figplane,draw=figgray!30,line width=.6pt]
   (-4.7,-4.7,0)--(4.7,-4.7,0)--(4.7,4.7,0)--(-4.7,4.7,0)--cycle;
 % Опорная круговая орбита лежит в плоскости OXY.
 \draw[draw=figgray!85,line width=.9pt,dash pattern=on 4pt off 2.5pt]
   plot[domain=0:360,samples=180,variable=\t]
   ({\orba*cos(\t)},{\orba*sin(\t)},0);
 \draw[guide] (-1,0,0)--(O);
 \draw[guide] (0,-.8,0)--(O);
 \draw[fixed] (O)--(5.8,0,0) node[below right=2pt] {$X$};
 \draw[fixed] (O)--(0,5.3,0) node[above left=2pt] {$Y$};
 \draw[fixed] (O)--(0,0,1.15) node[above=3pt] {$Z$};
 % Текущая орбита: нижняя полудуга находится под опорной плоскостью.
 \draw[draw=figorbit!75,line width=1pt,dash pattern=on 4pt off 2.5pt]
   plot[domain=180:360,samples=100,variable=\t]
   ({\orba*cos(\t)},{\orba*sin(\t)*cos(\orbi)},
    {\orba*sin(\t)*sin(\orbi)});
 \draw[draw=figorbit,line width=1.3pt]
   plot[domain=0:180,samples=100,variable=\t]
   ({\orba*cos(\t)},{\orba*sin(\t)*cos(\orbi)},
    {\orba*sin(\t)*sin(\orbi)});
 \draw[draw=figorbit,line width=1.3pt,-{Stealth[length=2.5mm]}]
   plot[domain=135:154,samples=20,variable=\t]
   ({\orba*cos(\t)},{\orba*sin(\t)*cos(\orbi)},
    {\orba*sin(\t)*sin(\orbi)});
 \coordinate (orbitlabel) at ({\orba*cos(115)},
   {\orba*sin(115)*cos(\orbi)},{\orba*sin(115)*sin(\orbi)});
 \coordinate (initiallabel) at ({\orba*cos(300)},{\orba*sin(300)},0);
 % r и u находятся в опорной плоскости; z — вдоль её нормали.
 \path[fill=figrad!10] (O)--(1.4,0,0)
   plot[domain=0:\phiu,samples=25,variable=\t]
   ({1.4*cos(\t)},{1.4*sin(\t)},0)--cycle;
 \draw[draw=figrad,line width=1pt] (O)--(H)
   node[pos=.63,above=4pt,text=figrad,lab] {$r$};
 \draw[guide] (H)--(P);
 \draw[draw=figink,line width=1.2pt,-{Stealth[length=2.6mm]}]
   (O)--(P) node[pos=.60,above left=3pt,lab] {$\boldsymbol R$};
 \draw[draw=figrad,line width=.9pt,-{Stealth[length=1.8mm]}]
   plot[domain=0:\phiu,samples=25,variable=\t]
   ({1.4*cos(\t)},{1.4*sin(\t)},0);
 \node[text=figrad] at ({1.83*cos(12)},{1.83*sin(12)},0) {$u$};
 % Фаза опорной точки в той же неподвижной плоскости.
 \draw[draw=figtan!80,line width=.8pt,dash pattern=on 3pt off 2pt] (O)--(Q);
 \draw[draw=figtan,line width=.9pt,-{Stealth[length=1.8mm]}]
   plot[domain=0:50,samples=40,variable=\t]
   ({2.35*cos(\t)},{2.35*sin(\t)},0);
 \node[text=figtan,lab] at ({2.68*cos(42)},{2.68*sin(42)},0) {$\varphi$};
 \node[point,fill=figtan] at (Q) {};
 \node[above left=3pt,text=figtan,lab] at (Q) {$Q$};
 % Положительные направления базиса, единичная длина в одном масштабе.
 \draw[local,figrad] (P)--++({1.8*cos(\phiu)},{1.8*sin(\phiu)},0)
   node[above right=4pt,align=left,text=figrad]
   {$\boldsymbol e_r$\\[2pt]$V_r=\dot r,\quad S$};
 \draw[local,figtan] (P)--++({-1.8*sin(\phiu)},{1.8*cos(\phiu)},0)
   node[above left=4pt,xshift=-3mm,text=figtan]
   {$\boldsymbol e_t,\quad V_t=r\dot u,\quad T$};
 \draw[local,fignormal] (P)--++(0,0,1.8)
   node[above=4pt,align=center,text=fignormal]
   {$\boldsymbol e_z$\\[2pt]$V_z=\dot z,\quad W$};
 % Вектор скорости направлен по касательной к текущей орбите.
 \draw[draw=figorbit,line width=1.55pt,-{Stealth[length=2.8mm,width=1.9mm]}]
   (P)--++({-1.8*sin(\orbv)},
   {1.8*cos(\orbv)*cos(\orbi)},{1.8*cos(\orbv)*sin(\orbi)})
   node[above left=3pt,text=figorbit,lab] {$\boldsymbol V$};
 \node[point] at (O) {};
 \node[below left=3pt] at (O) {$O$};
 \node[point,fill=figrad] at (H) {};
 \node[below right=4pt,text=figrad,lab] at (H) {$H$};
 \node[point,fill=figorbit] at (P) {};
 \node[right=4pt,xshift=4mm,yshift=-4mm,text=figorbit] at (P) {$P$};
\end{scope}

\coordinate (HZ) at ($(H)+(3.05,0)$);
\coordinate (PZ) at ($(P)+(3.05,0)$);
\draw[guide] ($(H)+(.14,0)$)--($(HZ)+(.2,0)$);
\draw[guide] ($(P)+(.14,0)$)--($(PZ)+(.2,0)$);
\draw[fignormal,line width=.85pt,{Stealth[length=1.8mm]}-{Stealth[length=1.8mm]}]
 (HZ)--(PZ) node[midway,right=4pt,text=fignormal] {$z$};
\draw[figgray,line width=.7pt]
 ($(H)+(-.23,-.040)$)--($(H)+(-.23,.18)$)--($(H)+(0,.22)$);

\node[anchor=west,text=figorbit,font=\small\bfseries] at (1.7,10.05)
 {Орбита КА};
\draw[draw=figorbit!70,line width=.6pt] (3.15,9.88)--(orbitlabel);
\node[anchor=west,text=figgray,align=left] at (13.15,2.9) {Опорная\\орбита};
\draw[draw=figgray!70,line width=.6pt] (13.05,2.9)--(initiallabel);
\node[text=figgray,font=\footnotesize] at (7.3,4.10) {Плоскость $OXY$};
\node[text=figink] at (7.6,.18)
 {$\boldsymbol R=r\boldsymbol e_r+z\boldsymbol e_z,\qquad
 s=\lVert\boldsymbol R\rVert=\sqrt{r^2+z^2}$};
\end{tikzpicture}
\endgroup
}
\end{column}
\begin{column}{.36\textwidth}
\RaggedRight
$OXY$ — неподвижная плоскость опорной орбиты. Ось $Z$ — нормаль к ней.\par\medskip
$P$ — КА, $H$ — его проекция на $OXY$; $Q$ — опорная точка на круге радиуса $r_0$.\par\medskip
Углы $u=\angle XOH$ и $\varphi=\angle XOQ$ отсчитываются от одной оси.\par\medskip
$r$ — расстояние до оси $Z$; $z$ — высота над опорной плоскостью.
\end{column}
\end{columns}
\end{frame}

\begin{frame}{Вращение базиса и кинематика}
\relax
\begin{columns}[T,onlytextwidth]
\begin{column}{.60\textwidth}
{\eqcompact\derivationmathspacing
\[
\begin{gathered}
 \boldsymbol e_r=(\cos u,\sin u,0),\qquad
 \boldsymbol e_t=(-\sin u,\cos u,0),\\[3pt]
 \boldsymbol e_z=(0,0,1),\\[20pt]
 \dot{\boldsymbol e}_r=\dot u\boldsymbol e_t,\qquad
 \dot{\boldsymbol e}_t=-\dot u\boldsymbol e_r,\qquad
 \dot{\boldsymbol e}_z=0,\\[20pt]
 \boldsymbol R=r\boldsymbol e_r+z\boldsymbol e_z,\\[5pt]
 \boldsymbol V=\dot r\boldsymbol e_r+r\dot u\boldsymbol e_t+\dot z\boldsymbol e_z,\\[20pt]
 \ddot{\boldsymbol R}=(\ddot r-r\dot u^2)\boldsymbol e_r
        +(r\ddot u+2\dot r\dot u)\boldsymbol e_t
        +\ddot z\boldsymbol e_z.
\end{gathered}
\]
}
\end{column}
\begin{column}{.37\textwidth}
\RaggedRight
При дифференцировании учитываем вращение базиса.\par\medskip
$\boldsymbol e_r$ направлен по проекции радиус-вектора на опорную плоскость.\par\medskip
$\boldsymbol e_z$ не обязан совпадать с нормалью к текущей орбите.
\end{column}
\end{columns}
\end{frame}

\begin{frame}{Точные уравнения движения}
\relax
\begin{columns}[T,onlytextwidth]
\begin{column}{.53\textwidth}
{\derivationmathspacing
\[
\begin{aligned}
 -\frac{\mu}{s^3}\boldsymbol R
 &=-\frac{\mu r}{s^3}\boldsymbol e_r-\frac{\mu z}{s^3}\boldsymbol e_z,\\[4pt]
 \boldsymbol a_{\text{упр}}&=S\boldsymbol e_r+T\boldsymbol e_t+W\boldsymbol e_z.
\end{aligned}
\]
\[
\boxed{\begin{aligned}
 \ddot r-r\dot u^2&=-\frac{\mu r}{(r^2+z^2)^{3/2}}+S,\\[7pt]
 r\ddot u+2\dot r\dot u&=T,\\[7pt]
 \ddot z&=-\frac{\mu z}{(r^2+z^2)^{3/2}}+W.
\end{aligned}}
\]
}
\end{column}
\begin{column}{.44\textwidth}
\RaggedRight
Компоненты управляющего ускорения:\par\smallskip
$S$ — радиальная;\par
$T$ — трансверсальная;\par
$W$ — вдоль $Z$.\par\medskip
У центральной силы нет проекции на $\boldsymbol e_t$.\par\medskip
При $S=T=W=0$ движение пассивное.
\end{column}
\end{columns}
\end{frame}

\begin{frame}{Координаты и физические скорости}
\relax
\begin{columns}[T,onlytextwidth]
\begin{column}{.48\textwidth}
{\eqcompact\derivationmathspacing
\[
\boxed{\begin{gathered}
 \dot r=V_r,\qquad \dot u=\frac{V_t}{r},\qquad \dot z=V_z,\\[20pt]
 \begin{aligned}
 \dot V_r&=\frac{V_t^2}{r}-\frac{\mu r}{(r^2+z^2)^{3/2}}+S,\\[6pt]
 \dot V_t&=-\frac{V_rV_t}{r}+T,\\[6pt]
 \dot V_z&=-\frac{\mu z}{(r^2+z^2)^{3/2}}+W.
 \end{aligned}
\end{gathered}}
\]
}
\end{column}
\begin{column}{.49\textwidth}
\RaggedRight
Проекции инерциальной скорости:
\[
 V_r=\dot r,\quad V_t=r\dot u,\quad V_z=\dot z.
\]
Дифференцируем $V_t=r\dot u$:
{\eqcompact\derivationmathspacing
\[
\begin{aligned}
 \dot V_t&=\dot r\dot u+r\ddot u\\[3pt]
 &=\dot r\dot u+(T-2\dot r\dot u)\\[3pt]
 &=T-\frac{V_rV_t}{r}.
\end{aligned}
\]
}
Здесь используются компоненты физической скорости в цилиндрическом базисе.
\end{column}
\end{columns}
\end{frame}

\begin{frame}{Опорный круг и отклонения}
\relax
\begin{columns}[T,onlytextwidth]
\begin{column}{.46\textwidth}
{\eqcompact\derivationmathspacing
\[
\begin{gathered}
 r_0=a_0=\frac{a_i+a_f}{2},\qquad V_0=\sqrt{\frac{\mu}{r_0}},\\[5pt]
 n_0=\sqrt{\frac{\mu}{r_0^3}}=\frac{V_0}{r_0},\\[5pt]
 \varphi(t)=n_0t+\varphi_0,\\[11pt]
 r=r_0+\Delta r,\\[5pt]
 V_r=\Delta V_r,\qquad V_t=V_0+\Delta V_t,\\[5pt]
 u=\varphi+\Delta u.
\end{gathered}
\]
}
\end{column}
\begin{column}{.51\textwidth}
\RaggedRight
В плоскости $OXY$ выбрана круговая опорная орбита. Индексы $i,f$ относятся к начальной и конечной орбитам.\par\medskip
$Q$ — условный аппарат, равномерно движущийся по опорному кругу.\par\medskip
$\Delta u=u-\varphi$ — угловое опережение опорной точки; при $\Delta u<0$ аппарат отстаёт.
\end{column}
\end{columns}
\end{frame}

\begin{frame}{Порядок приближения}
\relax
\begin{columns}[T,onlytextwidth]
\begin{column}{.56\textwidth}
{\derivationmathspacing
\[
\begin{gathered}
 \frac{|\Delta r|}{r_0},\ \frac{|\Delta V_r|}{V_0},\ \frac{|\Delta V_t|}{V_0}
       =O(\varepsilon),\\[7pt]
 \frac{|z|}{r_0},\ \frac{|V_z|}{V_0},\ |\Delta u|=O(\varepsilon),\\[7pt]
 S,T,W=O(\varepsilon n_0V_0),\qquad \varepsilon\ll1.
\end{gathered}
\]
\[
\begin{gathered}
 s=\sqrt{(r_0+\Delta r)^2+z^2}
   =r_0+\Delta r+O(\varepsilon^2r_0),\\[5pt]
 \frac{1}{r_0+\Delta r}
   \simeq\frac{1}{r_0}\left(1-\frac{\Delta r}{r_0}\right).
\end{gathered}
\]
}
\end{column}
\begin{column}{.41\textwidth}
\RaggedRight
Члены второго и более высоких порядков по $\varepsilon$ отбрасываются.\par\medskip
Управление также считается возмущением первого порядка на рассматриваемом интервале.\par\medskip
Малость координатных отклонений требуется на всём интервале движения.
\end{column}
\end{columns}
\end{frame}

\begin{frame}{Линеаризация правых частей}
\relax
\begin{columns}[T,onlytextwidth]
\begin{column}{.55\textwidth}
{\eqcompact\derivationmathspacing
\[
\begin{aligned}
 \frac{V_t^2}{r}
 &\simeq\frac{V_0^2}{r_0}+2n_0\Delta V_t-n_0^2\Delta r,\\[6pt]
 \frac{\mu r}{(r^2+z^2)^{3/2}}
 &\simeq\frac{\mu}{r_0^2}-2n_0^2\Delta r,\\[6pt]
 \frac{V_rV_t}{r}&\simeq n_0\Delta V_r,\\[6pt]
 \frac{\mu z}{(r^2+z^2)^{3/2}}&\simeq n_0^2z,\\[6pt]
 \frac{V_t}{r}&\simeq n_0+\frac{\Delta V_t-n_0\Delta r}{r_0}.
\end{aligned}
\]
}
\end{column}
\begin{column}{.42\textwidth}
\RaggedRight
Высота входит в радиальное ускорение через $z^2$: линейного члена по $z$ нет.\par\medskip
Постоянные члены сокращаются вследствие равенства
\[
 \frac{V_0^2}{r_0}=\frac{\mu}{r_0^2}.
\]
\end{column}
\end{columns}
\end{frame}

\begin{frame}{Линейная система отклонений}
\relax
\begin{columns}[T,onlytextwidth]
\begin{column}{.53\textwidth}
{\derivationmathspacing
\[
\boxed{\begin{aligned}
 \Delta\dot r&=\Delta V_r,\\[4pt]
 \Delta\dot V_r&=n_0^2\Delta r+2n_0\Delta V_t+S,\\[4pt]
 \Delta\dot V_t&=-n_0\Delta V_r+T,\\[4pt]
 \Delta\dot u&=\frac{\Delta V_t-n_0\Delta r}{r_0},\\[4pt]
 \dot z&=V_z,\\[4pt]
 \dot V_z&=-n_0^2z+W.
\end{aligned}}
\]
}
\end{column}
\begin{column}{.44\textwidth}
\RaggedRight
Система первого приближения.\par\medskip
Здесь и далее равенства для линеаризованной модели записываются без остаточных членов.\par\medskip
Их отсутствие не означает точности этих равенств для исходной нелинейной задачи.
\end{column}
\end{columns}
\end{frame}

\begin{frame}{Уравнения Клохесси--Уилтшира}
\relax
\begin{columns}[T,onlytextwidth]
\begin{column}{.56\textwidth}
{\eqcompact\derivationmathspacing
\[
\begin{aligned}
 \xi&=\Delta r,\qquad \eta=r_0\Delta u,\\[4pt]
 \dot\eta&=\Delta V_t-n_0\xi,\qquad
 \Delta V_t=\dot\eta+n_0\xi,\\[4pt]
 \ddot\xi&=n_0^2\xi+2n_0(\dot\eta+n_0\xi)+S,\\[4pt]
 \ddot\eta&=\Delta\dot V_t-n_0\dot\xi=-2n_0\dot\xi+T.
\end{aligned}
\]
\[
\boxed{\begin{aligned}
 \ddot\xi-2n_0\dot\eta-3n_0^2\xi&=S,\\[5pt]
 \ddot\eta+2n_0\dot\xi&=T,\\[5pt]
 \ddot z+n_0^2z&=W.
\end{aligned}}
\]
}
\end{column}
\begin{column}{.41\textwidth}
\RaggedRight
$\dot\eta$ — производная относительной продольной координаты во вращающейся системе.\par\medskip
$\Delta V_t$ — отклонение физической трансверсальной скорости.\par\medskip
Связь между ними:
\[
\dot\eta=\Delta V_t-n_0\Delta r.
\]
\end{column}
\end{columns}
\end{frame}

\begin{frame}{Нормировка пассивного движения}
\relax
\begin{columns}[T,onlytextwidth]
\begin{column}{.43\textwidth}
{\derivationmathspacing
\[
\boxed{\begin{aligned}
 (\drh)'&=\dvr,\\[5pt]
 (\dvr)'&=\drh+2\dvt,\\[5pt]
 (\dvt)'&=-\dvr,\\[5pt]
 (\Delta u)'&=\dvt-\drh,\\[5pt]
 (\zh)'&=\vz,\\[5pt]
 (\vz)'&=-\zh.
\end{aligned}}
\]
}
\end{column}
\begin{column}{.54\textwidth}
При $S=T=W=0$ вводим масштабы \mbox{$r_0$ и $V_0$}:
{\derivationmathspacing
\[
\begin{gathered}
 \drh=\frac{\Delta r}{r_0},\qquad \zh=\frac{z}{r_0},\\[5pt]
 \dvr=\frac{\Delta V_r}{V_0},\quad
 \dvt=\frac{\Delta V_t}{V_0},\quad
 \vz=\frac{V_z}{V_0}.
\end{gathered}
\]
\[
 \frac{\mathrm d}{\mathrm dt}=n_0\frac{\mathrm d}{\mathrm d\varphi}.
\]
}
\RaggedRight
Штрих обозначает дифференцирование по~$\varphi$. Крышка — нормировку координат на~$r_0$, скоростей на~$V_0$.
\end{column}
\end{columns}
\end{frame}

\begin{frame}{Энергия и большая полуось}
\relax
\begin{columns}[T,onlytextwidth]
\begin{column}{.58\textwidth}
{\eqcompact\derivationmathspacing
\[
\begin{aligned}
 \mathcal E&=\frac{\|\boldsymbol V\|^2}{2}-\frac{\mu}{s}
                =-\frac{\mu}{2a},\\[5pt]
 \mathcal E&\simeq-\frac{\mu}{2r_0}+V_0^2(\drh+\dvt),\\[5pt]
 -\frac{\mu}{2a}&\simeq-\frac{\mu}{2r_0}
             +\frac{V_0^2}{2}\frac{a-r_0}{r_0}.
\end{aligned}
\]
\[
 \boxed{\frac{a-r_0}{r_0}=2(\drh+\dvt).}
\]
\[
 \underbrace{(\drh+\dvt)'=0}_{\text{пассивный участок}}
 \quad\Longrightarrow\quad
 (\drh)''+\drh=\frac{a-r_0}{r_0}.
\]
}
\end{column}
\begin{column}{.39\textwidth}
\RaggedRight
$\mathcal E$ — удельная механическая энергия.\par\medskip
Приравнивая линейные части двух разложений, получаем связь с большой полуосью.\par\medskip
$a$ — размерная большая полуось текущей орбиты. На пассивном участке $a$ не меняется.
\end{column}
\end{columns}
\end{frame}

\begin{frame}{Решение на пассивном участке}
\relax
\begin{columns}[T,onlytextwidth]
\begin{column}{.59\textwidth}
{\eqcompact\derivationmathspacing
\[
\boxed{\begin{aligned}
 \drh&=\frac{a-r_0}{r_0}-e_x\cos\varphi-e_y\sin\varphi,\\[6pt]
 \dvr&=e_x\sin\varphi-e_y\cos\varphi,\\[6pt]
 \dvt&=e_x\cos\varphi+e_y\sin\varphi-\frac{a-r_0}{2r_0},\\[6pt]
 \zh&=i_x\sin\varphi-i_y\cos\varphi,\\[6pt]
 \vz&=i_x\cos\varphi+i_y\sin\varphi.
\end{aligned}}
\]
}
\end{column}
\begin{column}{.38\textwidth}
\RaggedRight
$a,e_x,e_y,i_x,i_y$ — пять постоянных интегрирования.\par\medskip
Шестая постоянная задаёт начальную фазу.\par\medskip
Геометрический смысл $e_x,e_y,i_x,i_y$ устанавливается через вектор эксцентриситета и нормаль к орбите.
\end{column}
\end{columns}
\end{frame}

\begin{frame}{Фаза на пассивном участке}
\relax
{\eqcompact\derivationmathspacing
\[
 (\Delta u)'=\dvt-\drh=-\frac32\frac{a-r_0}{r_0}
                       +2e_x\cos\varphi+2e_y\sin\varphi.
\]
\[
\boxed{\begin{aligned}
 \Delta u(\varphi)={}&\Delta u(\varphi_*)-
      \frac32\frac{a-r_0}{r_0}(\varphi-\varphi_*)\\[4pt]
 &+2e_x(\sin\varphi-\sin\varphi_*)
      -2e_y(\cos\varphi-\cos\varphi_*).
\end{aligned}}
\]
}
\begin{columns}[T,onlytextwidth]
\begin{column}{.49\textwidth}
\RaggedRight
$\varphi_*$ соответствует моменту $t_*$. Значение $\Delta u(\varphi_*)$ задаёт начальную фазу.\par\medskip
На пассивном участке $a,e_x,e_y$ постоянны.
\end{column}
\begin{column}{.48\textwidth}
\RaggedRight
Член, пропорциональный $\varphi-\varphi_*$, накапливает угловое расхождение.\par\medskip
Когда оно перестаёт быть малым, описание относительно одной опорной точки требует пересмотра.
\end{column}
\end{columns}
\end{frame}

\begin{frame}{Вектор эксцентриситета}
\relax
\begin{columns}[T,onlytextwidth]
\begin{column}{.58\textwidth}
{\eqcompact\derivationmathspacing
\[
\begin{aligned}
 \boldsymbol e&=\frac{\boldsymbol V\times(\boldsymbol R\times\boldsymbol V)}{\mu}
                 -\frac{\boldsymbol R}{s},\\[6pt]
 \|\boldsymbol V\|^2-\mu/s&\simeq V_0^2(\drh+2\dvt),\\[3pt]
 \boldsymbol R\cdot\boldsymbol V&\simeq r_0V_0\dvr,\\[6pt]
 \boldsymbol e&\simeq(\drh+2\dvt)\boldsymbol e_r-\dvr\boldsymbol e_t.
\end{aligned}
\]
\[
\boxed{\begin{aligned}
 e_x&=(\drh+2\dvt)\cos\varphi+\dvr\sin\varphi,\\[4pt]
 e_y&=(\drh+2\dvt)\sin\varphi-\dvr\cos\varphi.
\end{aligned}}
\]
}
\end{column}
\begin{column}{.39\textwidth}
\RaggedRight
$e_x,e_y$ — компоненты вектора эксцентриситета в неподвижных осях.\par\medskip
Замена $u$ на $\varphi$ при малых коэффициентах даёт ошибку второго порядка.\par\medskip
Для орбиты в $OXY$:
{\eqcompact\derivationmathspacing
\[
\begin{aligned}e_x&=e\cos\varpi,\\e_y&=e\sin\varpi.\end{aligned}
\]
}
$\varpi$ — долгота перицентра от оси $X$.
\end{column}
\end{columns}
\end{frame}

\begin{frame}{Нормаль и плоскость орбиты}
\relax
\begin{columns}[T,onlytextwidth]
\begin{column}{.53\textwidth}
{\eqcompact\derivationmathspacing
\[
\begin{aligned}
 \boldsymbol k&=\frac{\boldsymbol R\times\boldsymbol V}
                         {\|\boldsymbol R\times\boldsymbol V\|}\\[4pt]
 &\simeq\boldsymbol e_z-\zh\boldsymbol e_r-\vz\boldsymbol e_t
       =(i_y,-i_x,1),\\[7pt]
 i_x&=\zh\sin\varphi+\vz\cos\varphi,\\[4pt]
 i_y&=-\zh\cos\varphi+\vz\sin\varphi.
\end{aligned}
\]
\[
\boxed{\begin{aligned}
 \zh&=i_x\sin\varphi-i_y\cos\varphi,\\[4pt]
 \vz&=i_x\cos\varphi+i_y\sin\varphi.
\end{aligned}}
\]
}
\end{column}
\begin{column}{.44\textwidth}
\RaggedRight
$\boldsymbol k$ — единичная нормаль к текущей орбите.
{\eqcompact\derivationmathspacing
\[
\begin{aligned}i_x&=i\cos\Omega,\\i_y&=i\sin\Omega.\end{aligned}
\]
}
$i$ — малое наклонение к опорной плоскости; $\Omega$ — долгота восходящего узла.\par\medskip
Постоянные $i_x,i_y$ задают ориентацию орбиты на пассивном участке.
\end{column}
\end{columns}
\end{frame}

\begin{frame}{Производные орбитальных параметров}
\relax
\begin{columns}[T,onlytextwidth]
\begin{column}{.62\textwidth}
{\eqcompact\derivationmathspacing\centering
\noindent\(\displaystyle\begin{aligned}[t]
 & \dot a=2r_0\frac{\mathrm d}{\mathrm dt}(\drh+\dvt)
                =\frac{2r_0}{V_0}T,\\[3pt]
 & \frac{\mathrm d}{\mathrm dt}(\drh+2\dvt)
       =-n_0\dvr+\frac{2T}{V_0},\\[3pt]
 & \frac{\mathrm d}{\mathrm dt}\dvr
       =n_0(\drh+2\dvt)+\frac{S}{V_0},\\[3pt]
 & \dot e_x=\left[\frac{\mathrm d}{\mathrm dt}(\drh+2\dvt)+n_0\dvr\right]\cos\varphi\\[2pt]
 & \quad+\left[\frac{\mathrm d}{\mathrm dt}\dvr-n_0(\drh+2\dvt)\right]\sin\varphi\\[4pt]
 & =\frac{S\sin\varphi+2T\cos\varphi}{V_0}.
\end{aligned}\)\par
}
\end{column}
\begin{column}{.35\textwidth}
\RaggedRight
Дифференцируем связи $a,e_x,e_y$ с состоянием при $\dot\varphi=n_0$.\par\medskip
Члены свободного движения сокращаются. Остаются слагаемые с $S,T$.\par\medskip
Аналогично находим $\dot e_y$.
\end{column}
\end{columns}
\end{frame}

\begin{frame}{Полная система шести уравнений}
\relax
\begin{columns}[T,onlytextwidth]
\begin{column}{.61\textwidth}
{\eqcompact\derivationmathspacing
\[
\boxed{\begin{aligned}
 \dot a&=\frac{2r_0}{V_0}T,\\[4pt]
 \dot e_x&=\frac{S\sin\varphi+2T\cos\varphi}{V_0},\\[4pt]
 \dot e_y&=\frac{-S\cos\varphi+2T\sin\varphi}{V_0},\\[4pt]
 \dot z&=V_z,\\[4pt]
 \dot V_z&=-n_0^2z+W,\\[4pt]
 \dot u&=n_0-\frac{3n_0}{2}\frac{a-r_0}{r_0}
              +2n_0(e_x\cos\varphi+e_y\sin\varphi).
\end{aligned}}
\]
}
\end{column}
\begin{column}{.36\textwidth}
\RaggedRight
Шесть переменных: $a,e_x,e_y,z,V_z,u$.\par\medskip
Система описывает движение при заданном управлении.\par\medskip
На пассивном участке постоянны $a,e_x,e_y$; величины $z,V_z,u$ продолжают изменяться.
\end{column}
\end{columns}
\end{frame}

\begin{frame}{Импульс и отклонение скорости}
\relax
\begin{columns}[T,onlytextwidth]
\begin{column}{.53\textwidth}
{\eqcompact\derivationmathspacing
\[
\begin{aligned}
 \boldsymbol R(t_j+0)&=\boldsymbol R(t_j-0),\\[4pt]
 \boldsymbol V(t_j+0)&=\boldsymbol V(t_j-0)+\Delta\boldsymbol V_j^{\imp},\\[7pt]
 \Delta\widehat V_{rj}^{\imp}&=\Delta V_{rj}^{\imp}/V_0,\\[3pt]
 \Delta\widehat V_{tj}^{\imp}&=\Delta V_{tj}^{\imp}/V_0,\\[3pt]
 \Delta\widehat V_{zj}^{\imp}&=\Delta V_{zj}^{\imp}/V_0,\\[7pt]
 \Delta V_t(t_j+0)&=\Delta V_t(t_j-0)+\Delta V_{tj}^{\imp}.
\end{aligned}
\]
}
\end{column}
\begin{column}{.44\textwidth}
\RaggedRight
$j$ — номер манёвра. Положение непрерывно, скорость получает скачок.
\begin{block}{Отклонение и импульс}
$\Delta V_t=V_t-V_0$ — текущее отклонение трансверсальной скорости от круговой.
\par
$\Delta V_{tj}^{\imp}$ — сообщённый скачок этой скорости.
\end{block}
$\|\Delta\widehat{\boldsymbol V}_j^{\imp}\|=O(\varepsilon)$.
\end{column}
\end{columns}
\end{frame}

\begin{frame}{Скачки переменных состояния}
\relax
\begin{columns}[T,onlytextwidth]
\begin{column}{.62\textwidth}
{\eqcompact\derivationmathspacing
\[
\boxed{\begin{aligned}
 a_j^+-a_j^-&=2r_0\Delta\widehat V_{tj}^{\imp},\\[5pt]
 e_{xj}^+-e_{xj}^-&=\Delta\widehat V_{rj}^{\imp}\sin\varphi_j
                       +2\Delta\widehat V_{tj}^{\imp}\cos\varphi_j,\\[5pt]
 e_{yj}^+-e_{yj}^-&=-\Delta\widehat V_{rj}^{\imp}\cos\varphi_j
                       +2\Delta\widehat V_{tj}^{\imp}\sin\varphi_j,\\[5pt]
 z_j^+-z_j^-&=0,\\[5pt]
 V_{zj}^+-V_{zj}^-&=V_0\Delta\widehat V_{zj}^{\imp},\\[5pt]
 u_j^+-u_j^-&=0.
\end{aligned}}
\]
}
\end{column}
\begin{column}{.35\textwidth}
\RaggedRight
$a_j^-=a(t_j-0)$, $a_j^+=a(t_j+0)$; аналогично для остальных переменных.\par\medskip
$\varphi_j=\varphi(t_j)$ — фаза импульса.\par\medskip
Условия скачка подставлены в связи $a,e_x,e_y$ с состоянием.
\end{column}
\end{columns}
\end{frame}

\begin{frame}{Начальные условия отклика}
\relax
\begin{columns}[T,onlytextwidth]
\begin{column}{.51\textwidth}
{\eqcompact\derivationmathspacing
\[
\begin{aligned}
 \drh(t_j+0)&=0,\\[2pt]
 \dvr(t_j+0)&=\Delta\widehat V_{rj}^{\imp},\\[2pt]
 \dvt(t_j+0)&=\Delta\widehat V_{tj}^{\imp},\\[2pt]
 \zh(t_j+0)&=0,\\[2pt]
 \vz(t_j+0)&=\Delta\widehat V_{zj}^{\imp},\\[2pt]
 \Delta u(t_j+0)&=0.
\end{aligned}
\]
\[
 \theta=\varphi-\varphi_j=n_0(t-t_j).
\]
}
\RaggedRight
До момента $t_j$ аппарат движется по опорному кругу, причём $u(t_j)=\varphi(t_j)$. После импульса $S=T=W=0$.
\end{column}
\begin{column}{.46\textwidth}
\centering\resizebox{\linewidth}{!}{% Рисунок дословно из manual/near-circular-transfer.tex.
\begingroup
\definecolor{figink}{HTML}{25323F}
\definecolor{figgray}{HTML}{687582}
\definecolor{figrad}{HTML}{176AA6}
\definecolor{figtan}{HTML}{087F78}
\definecolor{fignormal}{HTML}{B57714}
\definecolor{figplane}{HTML}{EEF4F8}
\definecolor{figorbit}{HTML}{79539C}
\begin{tikzpicture}[x=1cm,y=1cm,>=Stealth,line cap=round,
 lab/.style={fill=white,inner sep=1pt}]
\path[use as bounding box] (-.55,-.45) rectangle (7.6,5.3);
\begin{scope}[shift={(.3,.25)}]
\draw[figgray!60] (0,0)--(6.2,0) node[below] {$X$};
\draw[figgray,dashed] plot[domain=0:102,samples=70] ({4.4*cos(\x)},{4.4*sin(\x)});
\coordinate (O) at (0,0);
\coordinate (Qj) at ({4.4*cos(15)},{4.4*sin(15)});
\coordinate (Q) at ({4.4*cos(63)},{4.4*sin(63)});
\coordinate (H) at ({4.15*cos(83)},{4.15*sin(83)});
\draw[figgray,densely dashed] (O)--(Qj);
\draw[figtan,line width=.8pt] (O)--(Q);
\draw[figrad,line width=.8pt] (O)--(H);
\draw[figrad,line width=.9pt,->] plot[domain=0:83,samples=50]
 ({.85*cos(\x)},{.85*sin(\x)});
\node[figrad,lab] at (1.05,.5) {$u$};
\draw[figtan,line width=.9pt,->] plot[domain=0:63,samples=50]
 ({1.65*cos(\x)},{1.65*sin(\x)});
\node[figtan,lab] at (1.85,.92) {$\varphi$};
\draw[figorbit,line width=1.1pt,->] plot[domain=15:63,samples=40]
 ({2.65*cos(\x)},{2.65*sin(\x)});
\node[figorbit,lab] at (2.38,1.93) {$\theta$};
\draw[figgray,line width=.7pt,->] plot[domain=63:83,samples=25]
 ({3.7*cos(\x)},{3.7*sin(\x)});
\node[figgray,lab] at (1.04,3.56) {$\Delta u$};
\fill[figgray] (Qj) circle(1.8pt);
\node[above right=2pt,align=left] at (Qj) {$Q_j$\\$\varphi_j$};
\fill[figtan] (Q) circle(1.8pt);
\node[above right=2pt,text=figtan] at (Q) {$Q(t)$};
\fill[figrad] (H) circle(1.8pt);
\node[above left=2pt,text=figrad] at (H) {$H(t)$};
\fill[figink] (O) circle(1.8pt);
\node[below left] at (O) {$O$};
\end{scope}
\end{tikzpicture}
\endgroup
}\par
\figcaption{$\theta$ — перемещение опорной точки после импульса; $\Delta u$ сравнивает КА и опорную точку в один момент.}
\end{column}
\end{columns}
\end{frame}

\begin{frame}{Вывод отклика на импульс}
\relax
\begin{columns}[T,onlytextwidth]
\begin{column}{.63\textwidth}
{\eqcompact\derivationmathspacing
\noindent\(\displaystyle\begin{aligned}[t]
 & (\drh+\dvt)'=0
     \quad\Longrightarrow\quad\drh+\dvt=\Delta\widehat V_{tj}^{\imp},\\[5pt]
 & (\drh)''+\drh=2\Delta\widehat V_{tj}^{\imp},\\[3pt]
 & \drh(0)=0,\qquad (\drh)'(0)=\Delta\widehat V_{rj}^{\imp},\\[5pt]
 & \drh=\Delta\widehat V_{rj}^{\imp}\sin\theta
                +2\Delta\widehat V_{tj}^{\imp}(1-\cos\theta),\\[6pt]
 & \dvr=(\drh)',\qquad
 \dvt=\Delta\widehat V_{tj}^{\imp}-\drh,\\[5pt]
 & (\Delta u)'=\Delta\widehat V_{tj}^{\imp}-2\drh,\\[5pt]
 & (\zh)''+\zh=0,\qquad
 \zh(0)=0,\quad(\zh)'(0)=\Delta\widehat V_{zj}^{\imp}.
\end{aligned}\)\par
}
\end{column}
\begin{column}{.34\textwidth}
\RaggedRight
Штрих — производная по $\theta$ или $\varphi$.\par\medskip
Из радиального решения восстанавливаем скорости.\par\medskip
Фазу интегрируем с $\Delta u(0)=0$. Нормальное движение находим независимо.
\end{column}
\end{columns}
\end{frame}

\begin{frame}{Полный отклик на один импульс}
\relax
\begin{columns}[T,onlytextwidth]
\begin{column}{.64\textwidth}
{\eqcompact\derivationmathspacing
\[
\boxed{\begin{aligned}
 \drh&=\Delta\widehat V_{rj}^{\imp}\sin\theta
                 +2\Delta\widehat V_{tj}^{\imp}(1-\cos\theta),\\[5pt]
 \dvr&=\Delta\widehat V_{rj}^{\imp}\cos\theta
                 +2\Delta\widehat V_{tj}^{\imp}\sin\theta,\\[5pt]
 \dvt&=-\Delta\widehat V_{rj}^{\imp}\sin\theta
                 +\Delta\widehat V_{tj}^{\imp}(2\cos\theta-1),\\[5pt]
 \zh&=\Delta\widehat V_{zj}^{\imp}\sin\theta,\\[5pt]
 \vz&=\Delta\widehat V_{zj}^{\imp}\cos\theta,\\[5pt]
 \Delta u&=2\Delta\widehat V_{rj}^{\imp}(\cos\theta-1)
                 +\Delta\widehat V_{tj}^{\imp}(4\sin\theta-3\theta).
\end{aligned}}
\]
}
\end{column}
\begin{column}{.33\textwidth}
\RaggedRight
$\theta=n_0(t-t_j)$.\par\medskip
Слева стоят отклонения фактического движения при старте с опорного круга.\par\medskip
При произвольном начальном состоянии это вклад одного импульса в полное решение.
\end{column}
\end{columns}
\end{frame}

\begin{frame}{Суперпозиция импульсных воздействий}
\relax
\begin{columns}[T,onlytextwidth]
\begin{column}{.55\textwidth}
{\eqcompact\derivationmathspacing
\[
\begin{aligned}
 \drh(t)&=\drh^{(0)}(t)+\sum_{j:\,t_j\le t}\drh^{(j)}(t),\\[7pt]
 \Delta u(t)&=\Delta u^{(0)}(t)+\sum_{j:\,t_j\le t}\Delta u^{(j)}(t),\\[8pt]
 t_i&\le t_1\le\ldots\le t_N\le t_f,\\[4pt]
 \varphi_i&=\varphi(t_i),\quad\varphi_j=\varphi(t_j),\quad
 \varphi_f=\varphi(t_f),\\[4pt]
 \theta_j&=\varphi_f-\varphi_j=n_0(t_f-t_j).
\end{aligned}
\]
}
\par\medskip\RaggedRight
Конечное состояние складывается из свободного движения и откликов на все импульсы.
\end{column}
\begin{column}{.42\textwidth}
\RaggedRight
$(0)$ — свободное движение; $(j)$ — отклик на импульс. До $t_j$ он равен нулю.\par\medskip
Опорный круг и фаза общие для всех слагаемых.\par\medskip
Состояния заданы до первого и после последнего импульсов.
\end{column}
\end{columns}
\end{frame}

\begin{frame}{Свободное движение до момента $t_f$}
\relax
{\eqcompact\derivationmathspacing
\[
\boxed{\begin{aligned}
 z_i^{\mathrm{free}}(t_f)
 &=z_i\cos(\varphi_f-\varphi_i)
       +\frac{V_{zi}}{n_0}\sin(\varphi_f-\varphi_i),\\[6pt]
 V_{zi}^{\mathrm{free}}(t_f)
 &=-n_0z_i\sin(\varphi_f-\varphi_i)
       +V_{zi}\cos(\varphi_f-\varphi_i),\\[6pt]
 u_i^{\mathrm{free}}(t_f)
 &=u_i+n_0(t_f-t_i)-\frac{3n_0}{2}\frac{a_i-r_0}{r_0}(t_f-t_i)\\[3pt]
 &\quad+2e_{xi}(\sin\varphi_f-\sin\varphi_i)
       -2e_{yi}(\cos\varphi_f-\cos\varphi_i).
\end{aligned}}
\]
}
\begin{columns}[T,onlytextwidth]
\begin{column}{.49\textwidth}
\RaggedRight
Начальное состояние: $a_i,e_{xi},e_{yi},z_i,V_{zi},u_i$. Параметры $a_i,e_{xi},e_{yi}$ сохраняются на пассивном участке.
\end{column}
\begin{column}{.48\textwidth}
\RaggedRight
При выборе опорной плоскости совпадающей с плоскостью начальной орбиты $z_i=V_{zi}=0$. Свободные нормальные компоненты тогда нулевые.
\end{column}
\end{columns}
\end{frame}

\begin{frame}{Шесть граничных условий}
\relax
{\fontsize{9}{11}\selectfont\derivationmathspacing\setlength{\jot}{1pt}
\[
\boxed{\begin{aligned}
 2\sum\nolimits_{j=1}^{N}\Delta\widehat V_{tj}^{\imp}
 &=\frac{a_f-a_i}{r_0},\\[2pt]
 \sum\nolimits_{j=1}^{N}\bigl(\Delta\widehat V_{rj}^{\imp}\sin\varphi_j
       +2\Delta\widehat V_{tj}^{\imp}\cos\varphi_j\bigr)
 &=e_{xf}-e_{xi},\\[2pt]
 \sum\nolimits_{j=1}^{N}\bigl(-\Delta\widehat V_{rj}^{\imp}\cos\varphi_j
       +2\Delta\widehat V_{tj}^{\imp}\sin\varphi_j\bigr)
 &=e_{yf}-e_{yi},\\[2pt]
 \sum\nolimits_{j=1}^{N}\Delta\widehat V_{zj}^{\imp}\sin\theta_j
 &=\frac{z_f-z_i^{\mathrm{free}}(t_f)}{r_0},\\[2pt]
 \sum\nolimits_{j=1}^{N}\Delta\widehat V_{zj}^{\imp}\cos\theta_j
 &=\frac{V_{zf}-V_{zi}^{\mathrm{free}}(t_f)}{V_0},\\[2pt]
 \sum\nolimits_{j=1}^{N}\bigl[2\Delta\widehat V_{rj}^{\imp}(\cos\theta_j-1)
       +\Delta\widehat V_{tj}^{\imp}(4\sin\theta_j-3\theta_j)\bigr]
 &=u_f-u_i^{\mathrm{free}}(t_f).
\end{aligned}}
\]
}
\RaggedRight
$\theta_j=n_0(t_f-t_j)$ — угол до конечного момента с учётом полного числа оборотов.
\end{frame}

\begin{frame}{Условия на конечную плоскость}
\relax
\begin{columns}[T,onlytextwidth]
\begin{column}{.55\textwidth}
{\eqcompact\derivationmathspacing
\[
\begin{aligned}
 z_f/r_0&=\Delta i\sin\varphi_f,\\[2pt]
 V_{zf}/V_0&=\Delta i\cos\varphi_f,\\[7pt]
 \sum\nolimits_j\Delta\widehat V_{zj}\sin(\varphi_f-\varphi_j)
     &=\Delta i\sin\varphi_f,\\[3pt]
 \sum\nolimits_j\Delta\widehat V_{zj}\cos(\varphi_f-\varphi_j)
     &=\Delta i\cos\varphi_f.
\end{aligned}
\]
\[
 \Updownarrow
\]
\[
\boxed{\begin{aligned}
 \sum\nolimits_j\Delta\widehat V_{zj}\sin\varphi_j&=0,\\[3pt]
 \sum\nolimits_j\Delta\widehat V_{zj}\cos\varphi_j&=\Delta i.
\end{aligned}}
\]
}
\end{column}
\begin{column}{.42\textwidth}
\RaggedRight
Начальная плоскость совпадает с опорной. Ось $X$ направлена к восходящему узлу конечной орбиты относительно начальной.\par\smallskip
$\Delta i\ge0$ — малый угол между плоскостями.\par\smallskip
Начиная с этого слайда опускаем «имп». Крышка означает деление импульса на $V_0$.
\end{column}
\end{columns}
\end{frame}

\begin{frame}{Пять условий перехода на орбиту}
\relax
\begin{columns}[T,onlytextwidth]
\begin{column}{.60\textwidth}
{\eqcompact\derivationmathspacing
\[
\boxed{\begin{aligned}
 \sum\nolimits_{j=1}^{N}2\Delta\widehat V_{tj}&=\Delta a,\\[7pt]
 \sum\nolimits_{j=1}^{N}(\Delta\widehat V_{rj}\sin\varphi_j
       +2\Delta\widehat V_{tj}\cos\varphi_j)&=\Delta e_x,\\[7pt]
 \sum\nolimits_{j=1}^{N}(-\Delta\widehat V_{rj}\cos\varphi_j
       +2\Delta\widehat V_{tj}\sin\varphi_j)&=\Delta e_y,\\[7pt]
 \sum\nolimits_{j=1}^{N}\Delta\widehat V_{zj}\sin\varphi_j&=0,\\[7pt]
 \sum\nolimits_{j=1}^{N}\Delta\widehat V_{zj}\cos\varphi_j&=\Delta i.
\end{aligned}}
\]
}
\end{column}
\begin{column}{.37\textwidth}
Различия орбит:
{\eqcompact\derivationmathspacing
\[
\begin{aligned}
 \Delta a&=\frac{a_f-a_i}{r_0},\\[3pt]
 \Delta e_x&=e_{xf}-e_{xi},\\[3pt]
 \Delta e_y&=e_{yf}-e_{yi}.
\end{aligned}
\]
}
\RaggedRight
$\Delta a$ безразмерно; $e_x,e_y$ — в общих осях.
\par\smallskip
Начало маневрирования и конечная фаза не назначены; время перехода свободно.
\end{column}
\end{columns}
\end{frame}

\begin{frame}{Конечная фаза и время перехода}
\relax
{\eqcompact\derivationmathspacing
\[
\boxed{\begin{aligned}
 u_f={}&u_i^{\mathrm{free}}(t_f)
    +\sum_{j=1}^{N}\bigl[2\Delta\widehat V_{rj}(\cos\theta_j-1)\\[2pt]
 &\quad+\Delta\widehat V_{tj}(4\sin\theta_j-3\theta_j)\bigr].
\end{aligned}}
\]
}
\begin{columns}[T,onlytextwidth]
\begin{column}{.49\textwidth}
\textcolor{courseaccent}{\textbf{Переход на орбиту}}\par\smallskip
\RaggedRight
Последняя строка полной системы определяет получившуюся фазу $u_f$, а не ограничивает управление.

\end{column}
\begin{column}{.48\textwidth}
\textcolor{courseaccent}{\textbf{Заданное конечное состояние}}\par\smallskip
\RaggedRight
Если конечное состояние задано, полная система даёт шесть граничных условий. Все правые части сравнивают состояния в один момент $t_f$.

\end{column}
\end{columns}
Все пять условий периодичны по $\varphi_j$ с периодом $2\pi$. В фазовом условии углы $\theta_j=n_0(t_f-t_j)$ содержат полное число витков.
\end{frame}

\begin{frame}{Оптимальное импульсное маневрирование}
\relax
При поиске оптимального импульсного перехода необходимо определить \cite{EdelbaumGobetzWashington1966}:
\medskip
\begin{enumerate}
 \setlength{\itemsep}{10pt}
 \item Количество импульсов.
 \item Точки их приложения.
 \item Направления импульсов скорости.
 \item Величины импульсов и распределение суммарной характеристической скорости между ними.
\end{enumerate}
\end{frame}

\begin{frame}{Верхняя оценка числа импульсов}
\relax
{\justifying
В линейной задаче без дополнительных ограничений для оптимального маневрирования достаточно числа импульсов, не превышающего числа независимых граничных условий \cite{Edelbaum1966,Neustadt1964,SternPotter1965}.\par}
\medskip
\begin{center}
\renewcommand{\arraystretch}{1.2}
\begin{tabular}{@{}l@{\qquad}l@{\qquad}c@{}}
\toprule
Задача & Параметры & Число импульсов\\
\midrule
Компланарный переход & $a,e_x,e_y$ & $\leq 3$\\
Пространственный переход & $a,e_x,e_y,i_x,i_y$ & $\leq 5$\\
Компланарная встреча & $a,e_x,e_y,u$ & $\leq 4$\\
Пространственная встреча & $a,e_x,e_y,i_x,i_y,u$ & $\leq 6$\\
\bottomrule
\end{tabular}
\end{center}
\smallskip
Для задачи встречи конечный момент времени задан.
\end{frame}

\begin{frame}{Число импульсов}
\relax
\begin{columns}[T,onlytextwidth]
\begin{column}{.48\textwidth}
\begin{block}{Один импульс}
\begin{itemize}
 \setlength{\itemsep}{5pt}
 \item Переход возможен только между орбитами, имеющими общую точку.
 \item Может уступать нескольким импульсам по затратам характеристической скорости.
 \item Одно включение ДУ упрощает выполнение манёвра и повышает надёжность.
\end{itemize}
\end{block}
\end{column}
\begin{column}{.48\textwidth}
\begin{block}{Несколько импульсов}
\begin{itemize}
 \setlength{\itemsep}{5pt}
 \item Позволяют перейти между орбитами, не имеющими общей точки.
 \item Могут уменьшить суммарную характеристическую скорость.
 \item Повторные включения ДУ предъявляют дополнительные требования к её ресурсу и надёжности.
\end{itemize}
\end{block}
\end{column}
\end{columns}
\end{frame}

\begin{frame}{Изменение большой полуоси}
\relax
\begin{itemize}
 \setlength{\itemsep}{7pt}
 \item Большую полуось изменяет только трансверсальная составляющая импульса.
 \item Величина изменения не зависит от точки приложения импульса.
 \item Радиальная и нормальная составляющие в первом приближении не изменяют большую полуось.
\end{itemize}
\medskip
Величина трансверсальной составляющей:
\[
 2\Delta\widehat V_t=\Delta a,
\]
\[
 \boxed{\Delta\widehat V_t=\frac12\Delta a.}
\]
\end{frame}

\begin{frame}{Изменение вектора эксцентриситета}
\relax
\begin{itemize}
 \setlength{\itemsep}{5pt}
 \item Вектор эксцентриситета изменяют трансверсальная и радиальная составляющие импульса.
 \item Изменение зависит от точки приложения импульса.
\end{itemize}
\[
 \Delta e=\sqrt{(\Delta e_x)^2+(\Delta e_y)^2}.
\]
\begin{columns}[T,onlytextwidth]
\begin{column}{.48\textwidth}
\begin{block}{Трансверсальный импульс}
\centering $\displaystyle |\Delta\widehat V_t|=\frac12\Delta e.$\par\smallskip
\RaggedRight
Одновременно изменяет большую полуось.
\end{block}
\end{column}
\begin{column}{.48\textwidth}
\begin{block}{Радиальный импульс}
\centering $\displaystyle \vphantom{\frac12}|\Delta\widehat V_r|=\Delta e.$\par\smallskip
\RaggedRight
Большая полуось не меняется в первом приближении.
\end{block}
\end{column}
\end{columns}
\end{frame}

\begin{frame}{Изменение эксцентриситета в апсидах}
\relax
{\justifying
Для изменения эксцентриситета при неизменном направлении перицентра трансверсальный импульс прикладывают в перицентре или апоцентре.\par}
\smallskip
\begin{columns}[T,onlytextwidth]
\begin{column}{.48\textwidth}
\begin{block}{В перицентре}
\[
 \Delta\widehat V_t=\frac{e_f-e_i}{2}.
\]
\begin{itemize}
 \setlength{\itemsep}{3pt}
 \item Положительный импульс увеличивает эксцентриситет.
 \item Отрицательный импульс уменьшает эксцентриситет.
\end{itemize}
\end{block}
\end{column}
\begin{column}{.48\textwidth}
\begin{block}{В апоцентре}
\[
 \Delta\widehat V_t=-\frac{e_f-e_i}{2}.
\]
\begin{itemize}
 \setlength{\itemsep}{3pt}
 \item Положительный импульс уменьшает эксцентриситет.
 \item Отрицательный импульс увеличивает эксцентриситет.
\end{itemize}
\end{block}
\end{column}
\end{columns}
\end{frame}

\begin{frame}{Поворот плоскости орбиты}
\relax
\begin{itemize}
 \setlength{\itemsep}{8pt}
 \item Плоскости орбит совмещаются одним нормальным импульсом на линии их пересечения.
 \item Модуль импульса определяется углом $\Delta i$ между плоскостями.
\end{itemize}
\medskip
\[
 \boxed{\Delta\widehat V_z=\pm\Delta i.}
\]
\medskip
\begin{itemize}
 \setlength{\itemsep}{8pt}
 \item Знак «$+$» соответствует восходящему узлу конечной орбиты относительно начальной, знак «$-$» — нисходящему.
 \item В первом приближении большая полуось и эксцентриситет не меняются.
\end{itemize}
\end{frame}

\begin{frame}{Радиальный импульс: орбита}
\relax
\begin{columns}[T,onlytextwidth]
\begin{column}{.50\textwidth}
{\justifying
Исходная орбита — круговая.\par}
\smallskip
\centering% Baranov, Sec. 1.4.1, Fig. 1.9; first-order orbit shape.
% Native size: 70 x 40 mm. Black: reference circle; blue dashed: perturbed orbit.
\begin{tikzpicture}[course diagram,x=1mm,y=1mm]
\path[use as bounding box] (-24,-20) rectangle (46,20);
\draw[vector line] (-18,0)--(30,0) node[above] {$X$};
\draw[vector line] (0,-17)--(0,18) node[left] {$Y$};
\draw[courseink,line width=.65pt] (0,0) circle[radius=13];
\draw[orbit line,dashed] plot[domain=0:360,samples=181,variable=\a]
 ({13*(1+.16*sin(\a))*cos(\a)},{13*(1+.16*sin(\a))*sin(\a)});
\draw[vector line] (0,0)--(130:13) node[midway,left] {$r_0$};
\draw[courseaccent,line width=1.1pt,->] (13,0)--(25,0)
 node[pos=.72,above=1mm] {$\Delta V_r$};
\draw[courseaccent,line width=1.1pt,->] (0,0)--(0,-6)
 node[pos=.7,right=1mm] {$\Delta\boldsymbol e$};
\draw[aux line] (0,15.08)--(16,15.08)--(19,17);
\node[anchor=west,inner sep=0pt] at (19,17) {$\Delta r_\alpha$};
\draw[aux line] (0,-10.92)--(17,-10.92)--(20,-14);
\node[anchor=west,inner sep=0pt] at (20,-14) {$\Delta r_\pi$};
\fill[courseaccent] (13,0) circle[radius=.65];
\node[anchor=north west,inner sep=1mm] at (13,0) {$P$};
\node[anchor=west,text=coursemuted,font=\fontsize{8}{10}\selectfont]
 at (-21,-19) {$\Delta r_\alpha=-\Delta r_\pi=r_0\Delta\widehat V_r>0$};
\end{tikzpicture}

\figcaption{Чёрная линия --- опорный круг; голубой пунктир --- орбита после импульса. Отклонения увеличены.}
\end{column}
\begin{column}{.47\textwidth}
Импульс $\Delta\widehat V_r>0$ приложен при $t=0$, $\theta=n_0t$.
{\eqcompact
\setlength{\abovedisplayskip}{3pt}\setlength{\belowdisplayskip}{3pt}
\setlength{\abovedisplayshortskip}{2pt}\setlength{\belowdisplayshortskip}{2pt}
\[
 \begin{aligned}
 \Delta r&=r_0{\Delta\widehat V_r}\sin\theta,\\
 \Delta a&=0,\\
 \Delta\boldsymbol e&=-{\Delta\widehat V_r}\boldsymbol e_t.
 \end{aligned}
\]
}
Перицентр расположен за $90^\circ$ до точки приложения импульса.
{\eqcompact
\setlength{\abovedisplayskip}{3pt}\setlength{\belowdisplayskip}{3pt}
\setlength{\abovedisplayshortskip}{2pt}\setlength{\belowdisplayshortskip}{2pt}
\[
 \Delta r_\alpha=r_0{\Delta\widehat V_r},\qquad \Delta r_\pi=-r_0{\Delta\widehat V_r}.
\]
}
{\justifying
Здесь $\alpha$ --- апоцентр, $\pi$ --- перицентр.\par}
\end{column}
\end{columns}
\end{frame}

\begin{frame}{Радиальный импульс: фаза}
\relax
\begin{columns}[T,onlytextwidth]
\begin{column}{.50\textwidth}
\centering% Baranov, Sec. 1.4.1, Fig. 1.10. Native size: 70 x 40 mm.
\begin{tikzpicture}[course diagram,x=1mm,y=1mm]
\path[use as bounding box] (-9,-28) rectangle (61,12);
\draw[vector line] (0,-24)--(0,9) node[above right] {$\Delta u$};
\draw[vector line] (0,0)--(57,0) node[above] {$\theta$};
\draw[aux line,dashed] (25,0)--(25,-20);
\draw[orbit line] plot[domain=0:360,samples=121,variable=\a]
 ({\a/7.2},{10*(cos(\a)-1)});
\foreach \x/\lab in {25/\pi,50/2\pi}{
 \draw[courseink,line width=.5pt] (\x,-.8)--(\x,.8);
 \node[above=1mm,inner sep=0pt] at (\x,0) {$\lab$};}
\node[below=1mm,inner sep=0pt] at (25,-20) {$-4\Delta\widehat V_r$};
\fill[courseaccent] (25,-20) circle[radius=.65];
\node[below left,inner sep=1mm] at (0,0) {$0$};
\end{tikzpicture}

\end{column}
\begin{column}{.47\textwidth}
Для $\Delta\widehat V_r>0$:
{\eqcompact
\setlength{\abovedisplayskip}{3pt}\setlength{\belowdisplayskip}{3pt}
\setlength{\abovedisplayshortskip}{2pt}\setlength{\belowdisplayshortskip}{2pt}
\[
 \Delta u=2{\Delta\widehat V_r}(\cos\theta-1).
\]
}
Максимальное отставание при $\theta=\pi$:
{\eqcompact
\setlength{\abovedisplayskip}{3pt}\setlength{\belowdisplayskip}{3pt}
\setlength{\abovedisplayshortskip}{2pt}\setlength{\belowdisplayshortskip}{2pt}
\[
 \Delta u_{\min}=-4\Delta\widehat V_r.
\]
}
Через один опорный период:
{\eqcompact
\setlength{\abovedisplayskip}{3pt}\setlength{\belowdisplayskip}{3pt}
\setlength{\abovedisplayshortskip}{2pt}\setlength{\belowdisplayshortskip}{2pt}
\[
 \Delta u(2\pi)=0,\qquad T_0=\frac{2\pi}{n_0}.
\]
}
{\justifying
Вековое смещение в первом приближении отсутствует.\par}
\end{column}
\end{columns}
\end{frame}

\begin{frame}{Радиальный импульс: траектория}
\relax
Движение относительно КА на опорном круге, $\eta=r_0\Delta u$.
\smallskip
\begin{columns}[T,onlytextwidth]
\begin{column}{.50\textwidth}
\centering% Baranov, Sec. 1.4.1, Fig. 1.11(a). Native size: 70 x 40 mm.
% xi = eta/(r0 q), rho = Delta r/(r0 q), q = Delta V_r/V0 > 0.
\begin{tikzpicture}[course diagram,x=1mm,y=1mm]
\path[use as bounding box] (-59,-20) rectangle (11,20);
\draw[vector line] (-54,0)--(8,0) node[above] {$\eta$};
\draw[vector line] (0,-16)--(0,17) node[above right] {$\Delta r$};
\draw[orbit line] (-24,0) ellipse[x radius=24,y radius=12];
\draw[courseaccent,line width=1pt,->] plot[domain=43:59,samples=12,variable=\a]
 ({24*(cos(\a)-1)},{12*sin(\a)});
\draw[courseaccent,line width=1pt,->] plot[domain=223:239,samples=12,variable=\a]
 ({24*(cos(\a)-1)},{12*sin(\a)});
\draw[aux line,dashed] (-24,-12)--(-24,12);
\fill[courseaccent] (0,0) circle[radius=.7];
\node[anchor=north west,inner sep=1mm] at (0,0) {$0,\ T_0$};
\node[above=1mm,inner sep=0pt] at (-24,12) {$T_0/4$};
\node[below=1mm,inner sep=0pt] at (-24,-12) {$3T_0/4$};
\node[below=1.2mm,inner sep=0pt] at (-48,0) {$T_0/2$};
\draw[aux line] (-48,0)--(-48,14);
\node[above=1mm,inner sep=0pt] at (-46,14) {$-4r_0\Delta\widehat V_r$};
\node[anchor=south west,inner sep=0pt] at (-18,12) {$r_0\Delta\widehat V_r$};
\node[anchor=north west,inner sep=0pt] at (-18,-12) {$-r_0\Delta\widehat V_r$};
\end{tikzpicture}

\end{column}
\begin{column}{.47\textwidth}
{\eqcompact
\setlength{\abovedisplayskip}{3pt}\setlength{\belowdisplayskip}{3pt}
\setlength{\abovedisplayshortskip}{2pt}\setlength{\belowdisplayshortskip}{2pt}
\[
 \left(\frac{\eta+2r_0{\Delta\widehat V_r}}{2r_0{\Delta\widehat V_r}}\right)^2
 +\left(\frac{\Delta r}{r_0{\Delta\widehat V_r}}\right)^2=1.
\]
}
\begin{itemize}
 \setlength{\itemsep}{2pt}
 \item $T_0/4$: наибольшее увеличение радиуса.
 \item $T_0/2$: наибольшее отставание, $\Delta r=0$.
 \item $3T_0/4$: наименьший радиус.
 \item $T_0$: возврат в исходную относительную точку.
\end{itemize}
\end{column}
\end{columns}
\end{frame}

\begin{frame}{Сближение радиальными импульсами}
\relax
\begin{columns}[T,onlytextwidth]
\begin{column}{.50\textwidth}
\centering% Baranov, Sec. 1.4.1, Fig. 1.11(b). Native size: 70 x 40 mm.
% Repeated negative radial impulses: lower half-ellipses A -> B -> C -> O.
% With no impulse at C, passive upper half-ellipse returns to B.
% Equal arcs: Delta V_r < 0 at A and O; 2 Delta V_r at B and C.
% All axes have the same scale; the radial/along-track semiaxis ratio is 1:2.
\begin{tikzpicture}[course diagram,x=1mm,y=1mm]
\path[use as bounding box] (-63,-21) rectangle (7,19);
\draw[vector line] (-61,0)--(6,0) node[above] {$\eta$};
\draw[vector line] (0,-15)--(0,17) node[above right] {$\Delta r$};
\foreach \c in {-47.5,-28.5,-9.5}{
 \draw[orbit line] plot[domain=180:360,samples=61,variable=\a]
  ({\c+9.5*cos(\a)},{4.75*sin(\a)});
 \draw[courseaccent,line width=1pt,->] plot[domain=258:276,samples=9,variable=\a]
  ({\c+9.5*cos(\a)},{4.75*sin(\a)});}
\draw[courseink,line width=.7pt,dashed] plot[domain=0:180,samples=61,variable=\a]
 ({-28.5+9.5*cos(\a)},{4.75*sin(\a)});
\draw[courseink,line width=.7pt,->] plot[domain=80:98,samples=9,variable=\a]
 ({-28.5+9.5*cos(\a)},{4.75*sin(\a)});
\foreach \x/\lab in {-57/A,-38/B,0/O}{
 \fill[courseaccent] (\x,0) circle[radius=.7];
 \node[above left,inner sep=1mm] at (\x,0) {$\lab$};}
\fill[courseaccent] (-19,0) circle[radius=.7];
\node[above right,inner sep=1mm] at (-19,0) {$C$};
\foreach \x/\h/\lab in {-57/4/{\Delta V_r},-38/8/{2\Delta V_r},-19/8/{2\Delta V_r},0/4/{\Delta V_r}}{
 \draw[courseaccent,line width=.9pt,->] (\x,0)--(\x,-\h);
 \ifnum\x=0
  \node[below right=1mm,inner sep=0pt,text=courseaccent] at (\x,-\h) {$\lab$};
 \else
  \node[below=1mm,inner sep=0pt,text=courseaccent] at (\x,-\h) {$\lab$};
 \fi}
\node[text=coursemuted,font=\fontsize{8}{10}\selectfont]
 at (-29,15) {без импульса в $C$: возврат в $B$};
\node[text=coursemuted,font=\fontsize{8}{10}\selectfont]
 at (-28,-17) {Радиальные импульсы, $\Delta V_r<0$};
\end{tikzpicture}

\end{column}
\begin{column}{.47\textwidth}
{\justifying
В точке $A$ КА находится на орбите цели, позади неё. Импульс $\Delta V_r<0$ начинает движение по нижней половине эллипса.\par}
\smallskip
Через полвитка КА возвращается на исходный радиус и приближается к цели на $4r_0|\Delta\widehat V_r|$.
\par\smallskip
{\justifying
В $B,C$ импульсы $2\Delta V_r$ меняют знак радиальной скорости. В $O$ импульс $\Delta V_r$ гасит её.\par}
\end{column}
\end{columns}
\smallskip
{\justifying
\textbf{При пропуске импульса в $C$} КА возвращается в $B$ по верхней половине эллипса.\par}
\end{frame}

\begin{frame}{Трансверсальный импульс: орбита}
\relax
\begin{columns}[T,onlytextwidth]
\begin{column}{.50\textwidth}
\centering% Baranov, Sec. 1.4.2, Fig. 1.12; first-order orbit shape.
% Native size: 70 x 40 mm. Black: reference circle; blue dashed: perturbed orbit.
\begin{tikzpicture}[course diagram,x=1mm,y=1mm]
\path[use as bounding box] (-33,-20) rectangle (37,20);
\draw[vector line] (-25,0)--(26,0) node[above] {$X$};
\draw[vector line] (0,-18)--(0,18) node[left] {$Y$};
\draw[courseink,line width=.65pt] (0,0) circle[radius=12];
\draw[orbit line,dashed] plot[domain=0:360,samples=181,variable=\a]
 ({12*(1+.32*(1-cos(\a)))*cos(\a)},
  {12*(1+.32*(1-cos(\a)))*sin(\a)});
\draw[vector line] (0,0)--(125:12) node[midway,left] {$r_0$};
\draw[courseaccent,line width=1.1pt,->] (12,0)--(12,11)
 node[pos=.75,right=1mm] {$\Delta V_t$};
\draw[courseaccent,line width=1.1pt,->] (0,0)--(7,0)
 node[midway,below=1mm] {$\Delta\boldsymbol e$};
\draw[aux line] (-19.68,0)--(-23,8)--(-28,8);
\node[anchor=south,inner sep=1mm] at (-25,8) {$\Delta r_\alpha$};
\fill[courseaccent] (12,0) circle[radius=.65];
\node[anchor=north west,inner sep=1mm] at (12,0) {$P$};
\node[anchor=west,text=coursemuted,font=\fontsize{8}{10}\selectfont]
 at (-26,-19) {$\Delta r_\alpha=4r_0\Delta\widehat V_t,\quad\Delta r_\pi=0$};
\end{tikzpicture}

\figcaption{Чёрная линия --- опорный круг; голубой пунктир --- орбита после импульса. Отклонения увеличены.}
\end{column}
\begin{column}{.47\textwidth}
Положительный трансверсальный импульс: $\Delta\widehat V_t>0$.
{\eqcompact
\setlength{\abovedisplayskip}{3pt}\setlength{\belowdisplayskip}{3pt}
\setlength{\abovedisplayshortskip}{2pt}\setlength{\belowdisplayshortskip}{2pt}
\[
 \begin{aligned}
 \Delta r&=2r_0{\Delta\widehat V_t}(1-\cos\theta),\\
 \Delta a&=2{\Delta\widehat V_t}.
 \end{aligned}
\]
}
{\justifying
Точка импульса становится перицентром новой орбиты; противоположная точка --- апоцентром.\par}
{\eqcompact
\setlength{\abovedisplayskip}{3pt}\setlength{\belowdisplayskip}{3pt}
\setlength{\abovedisplayshortskip}{2pt}\setlength{\belowdisplayshortskip}{2pt}
\[
 \Delta r_\pi=0,\qquad \Delta r_\alpha=4r_0{\Delta\widehat V_t}.
\]
}
{\justifying
При ${\Delta\widehat V_t}<0$ точка импульса становится апоцентром; радиус противоположной апсиды уменьшается.\par}
\end{column}
\end{columns}
\end{frame}

\begin{frame}{Трансверсальный импульс: фаза}
\relax
\begin{columns}[T,onlytextwidth]
\begin{column}{.50\textwidth}
\centering% Baranov, Sec. 1.4.2, Fig. 1.13. Native size: 70 x 40 mm.
\begin{tikzpicture}[course diagram,x=1mm,y=1mm]
\path[use as bounding box] (-9,-28) rectangle (61,12);
\draw[vector line] (0,-25)--(0,9) node[above right] {$\Delta u$};
\draw[vector line] (0,0)--(57,0) node[above] {$\theta$};
\draw[coursemuted,line width=.55pt,dashed] (0,0)--(50,-21);
\draw[aux line,dashed] (50,0)--(50,-21);
\draw[orbit line] plot[domain=0:360,samples=181,variable=\a]
 ({\a/7.2},{21/(6*pi)*(4*sin(\a)-3*\a*pi/180)});
\foreach \x/\lab in {25/\pi,50/2\pi}{
 \draw[courseink,line width=.5pt] (\x,-.8)--(\x,.8);
 \node[above=1mm,inner sep=0pt] at (\x,0) {$\lab$};}
\node[below=1mm,inner sep=0pt] at (40,-21) {$-6\pi\Delta\widehat V_t$};
\draw[aux line] (34,-9)--(34,-14.28);
\node[above,inner sep=1mm,text=coursemuted] at (34,-9)
 {$-3\theta\Delta\widehat V_t$};
\fill[courseaccent] (50,-21) circle[radius=.65];
\node[below left,inner sep=1mm] at (0,0) {$0$};
\end{tikzpicture}

\end{column}
\begin{column}{.47\textwidth}
Для $\Delta\widehat V_t>0$:
{\eqcompact
\setlength{\abovedisplayskip}{3pt}\setlength{\belowdisplayskip}{3pt}
\setlength{\abovedisplayshortskip}{2pt}\setlength{\belowdisplayshortskip}{2pt}
\[
 \Delta u={\Delta\widehat V_t}(4\sin\theta-3\theta).
\]
}
Максимум опережения:
{\eqcompact
\setlength{\abovedisplayskip}{3pt}\setlength{\belowdisplayskip}{3pt}
\setlength{\abovedisplayshortskip}{2pt}\setlength{\belowdisplayshortskip}{2pt}
\[
 \begin{aligned}
 \theta_{\max}&\approx41^\circ24'35'',\\[3pt]
 \Delta u_{\max}&\approx0{,}4776\Delta\widehat V_t.
 \end{aligned}
\]
}
После $\theta\approx73^\circ05'32''$ начинается отставание.
\smallskip

За один опорный период:
{\eqcompact
\setlength{\abovedisplayskip}{3pt}\setlength{\belowdisplayskip}{3pt}
\setlength{\abovedisplayshortskip}{2pt}\setlength{\belowdisplayshortskip}{2pt}
\[
 \Delta u(2\pi)=-6\pi {\Delta\widehat V_t}.
\]
}
\end{column}
\end{columns}
\end{frame}

\begin{frame}{Трансверсальный импульс: траектория}
\relax
\centering% Baranov, Sec. 1.4.2, Fig. 1.14. Native size: 140 x 36 mm.
% r0 = 6700 km, Delta V_t = 0.002 km/s, mu = 398600.4418 km^3/s^2.
% eta = (Delta V_t/n0)(4 sin theta - 3 theta),
% Delta r = 2(Delta V_t/n0)(1-cos theta); axes both in kilometres.
\begin{tikzpicture}[course diagram,x=1mm,y=1mm]
\path[use as bounding box] (-122,-8) rectangle (18,28);
\pgfmathsetmacro{\impulselength}{1.7372939683809365}
\draw[vector line] (-119,0)--(15,0) node[above] {$\eta$, км};
\draw[vector line] (0,-5)--(0,25) node[above left] {$\Delta r$, км};
\foreach \a in {-40,-30,-20,-10}{
 \draw[courseink,line width=.5pt] ({2.6*\a},-.8)--({2.6*\a},.8);
 \node[below=1.2mm,inner sep=0pt] at ({2.6*\a},0) {$\a$};}
\foreach \a in {3.5,7}{
 \draw[courseink,line width=.5pt] (-.8,{3*\a})--(.8,{3*\a});
 \node[right=1.2mm,inner sep=0pt] at (0,{3*\a}) {$\a$};}
\begin{scope}
\clip (-119,-.3) rectangle (6,24);
\draw[orbit line] plot[domain=0:360,samples=181,variable=\a]
 ({2.6*\impulselength*(4*sin(\a)-3*\a*pi/180)},
  {6*\impulselength*(1-cos(\a))});
\draw[courseaccent,line width=.9pt,dashed] plot[domain=360:550,samples=101,variable=\a]
 ({2.6*\impulselength*(4*sin(\a)-3*\a*pi/180)},
  {6*\impulselength*(1-cos(\a))});
\draw[courseaccent,line width=1pt,->] plot[domain=114:122,samples=9,variable=\a]
 ({2.6*\impulselength*(4*sin(\a)-3*\a*pi/180)},
  {6*\impulselength*(1-cos(\a))});
\end{scope}
\fill[courseaccent] (0,0) circle[radius=.65];
\node[below right,inner sep=1mm] at (0,0) {$0$};
\fill[courseaccent] ({-2.6*\impulselength*6*pi},0) circle[radius=.65];
\node[below=1.2mm,inner sep=0pt] at ({-2.6*\impulselength*6*pi},0) {$T_0$};
\node[text=coursemuted,font=\fontsize{8}{10}\selectfont] at (-103,25)
 {следующий виток};
\end{tikzpicture}
\par
\raggedright
$\Delta V_t=2$ м/с, $r_0=6700$ км.
\smallskip

$\Delta r_{\max}\approx6{,}95$ км;
$\eta(T_0)\approx-32{,}75$ км.
\smallskip

{\justifying
Сначала КА смещается вперёд и наружу, затем отстаёт, продолжая подъём. После апоцентра он снижается до исходного радиуса позади невозмущённого КА. Пунктир --- начало следующего витка.\par}
\end{frame}

\begin{frame}{Нормальный импульс}
\relax
\begin{columns}[T,onlytextwidth]
\begin{column}{.50\textwidth}
\centering% Baranov, Sec. 1.4.3, Fig. 1.16. Native size: 70 x 40 mm.
% Oblique projection; the line of nodes is the x axis.
% The change of radius is second order and is omitted.
\begin{tikzpicture}[course diagram,x=1mm,y=1mm]
\path[use as bounding box] (-32,-20) rectangle (38,20);
\draw[vector line] (0,0)--(-14,-16) node[below left] {$X$};
\draw[vector line] (0,0)--(30,0) node[right] {$Y$};
\draw[vector line] (0,0)--(0,18) node[above left] {$Z$};
\draw[courseink,line width=.65pt] plot[domain=0:360,samples=181,variable=\a]
 ({-10*cos(\a)+23*sin(\a)},{-11*cos(\a)});
\draw[orbit line,dashed] plot[domain=0:360,samples=181,variable=\a]
 ({-10*cos(\a)+23*sin(\a)},{-11*cos(\a)+7*sin(\a)});
\draw[courseaccent,line width=1.1pt,->] (-10,-11)--(-10,-2.5)
 node[pos=.8,left=1mm] {$\Delta V_z$};
\fill[courseaccent] (-10,-11) circle[radius=.7];
\node[below right,inner sep=1mm] at (-10,-11) {$P$};
\draw[aux line] (0,0)--(23,7);
\draw[courseaccent,line width=.7pt] (10,0) arc[start angle=0,end angle=16.93,radius=10];
\node[inner sep=0pt] at (13.1,1.8) {$\Delta i$};
\draw[courseaccent,line width=.65pt,<->] (24.5,0)--(24.5,7)
 node[midway,right=1mm] {$|z|_{\max}$};
\draw[aux line] (23,7)--(26,7);
\node[text=coursemuted,font=\fontsize{8}{10}\selectfont] at (3,-18)
 {$|z|_{\max}=r_0|\Delta\widehat V_z|$};
\end{tikzpicture}

\figcaption{Чёрная линия --- опорный круг; голубой пунктир --- орбита после импульса. Угол поворота увеличен.}
\end{column}
\begin{column}{.47\textwidth}
После нормального импульса $\Delta\widehat V_z$:
{\eqcompact
\setlength{\abovedisplayskip}{3pt}\setlength{\belowdisplayskip}{3pt}
\[
 \begin{aligned}
 z&=r_0{\Delta\widehat V_z}\sin\theta,\\
 V_z&=V_0{\Delta\widehat V_z}\cos\theta.
 \end{aligned}
\]
}
Угол поворота плоскости: $\Delta i=|\Delta\widehat V_z|$.
\smallskip

Амплитуда отклонения: $|z|_{\max}=r_0|{\Delta\widehat V_z}|$.
\smallskip

{\justifying
В первом порядке $\Delta r=\Delta u=0$: движение относительно опорного КА состоит из колебаний по нормали.\par}

\end{column}
\end{columns}
\end{frame}

\begin{frame}{Источники}
\setbeamertemplate{bibliography item}{\insertbiblabel}
\begin{thebibliography}{9}
\bibitem{Baranov}
Баранов А. А. Маневрирование космических аппаратов в окрестности круговой орбиты // М.: Спутник. – 2016.
\bibitem{EdelbaumGobetzWashington1966}
Edelbaum T. N., Gobetz F. W., Washington M.
\href{https://ntrs.nasa.gov/citations/19660030604}{Minimum-impulse time-free transfer between elliptic orbits}.
NASA CR-636, 1966.
\bibitem{Edelbaum1966}
Edelbaum T. N.
\href{https://doi.org/10.2514/6.1966-7}{How many impulses?}
AIAA Paper 66-7, 1966.
\bibitem{Neustadt1964}
Neustadt L. W.
\href{https://doi.org/10.1137/0302004}{Optimization, a moment problem, and nonlinear programming}.
Journal of the Society for Industrial and Applied Mathematics, Series A: Control.
1964. Vol.~2, No.~1. P.~33--53.
\bibitem{SternPotter1965}
Stern R. G., Potter J. E.
\href{https://www.sciencedirect.com/science/article/abs/pii/S1474667017690700}{Optimization of midcourse velocity corrections}.
IFAC Proceedings Volumes. 1965. Vol.~2, No.~1. P.~70--83.
\end{thebibliography}
\end{frame}

\end{document}
